% ===== HAKKINDA =====
\chapter*{Hakkında}
\addcontentsline{toc}{chapter}{Hakkında}
\markboth{Hakkında}{Hakkında}

Ben Dost Seferoğlu. TÜBİTAK Fen Lisesi'nin ilk mezunlarındanım ve Koç Üniversitesi Elektrik-Elektronik Mühendisliği Bölümü'nde lisans eğitimime başlıyorum.

Ortaokulda yazılımla hiçbir ilgim yoktu; bilgisayar olimpiyatlarına lise hazırlık sınıfının sonunda başladım. Okulun ilk öğrencileri olduğumuz için önümüzde bize yol gösterecek bir üst dönem yoktu. İlk bir buçuk yılımı tek başıma, ne kadar ilerlediğimi kestiremeden ve gereksiz yere yüzlerce kolay soru çözerek geçirdim. Doğru bir yönlendirme alamadığım için de ilk senemde 1. Aşama Sınavı'nı geçemedim.

Benim için asıl dönüm noktası, alt dönemlerin okula gelişiyle başladı. Yurdun ortak alanında tek başıma problem çözerken yanıma geldiler ve birlikte çalışmaya başladık. Ders çıkışlarında ve gece etütlerinde saatlerce süren bu ortak çalışma, zamanla okulumuzda bir kültüre dönüştü; birbirimize bilmediklerimizi anlatıp eksiklerimizi kapatacak bir ortam kurduk. Bir süre sonra okuldaki alt dönemlerime ders vermeye, lisenin son yılında ise TÜBİTAK kamplarında eğitmenlik yapmaya başladım. Bu süreçte IOI, Balkan ve Asya-Pasifik Bilgisayar Olimpiyatları gibi yarışmalarda ülkemizi temsil etme fırsatı buldum.

Kendimi hiçbir zaman bu işi kolayca kavrayan o "parlak" öğrencilerden biri olarak görmedim; aksine çoğu kavramı zihnimde oturtabilmek için çok uğraşmam gerekti. Ancak bir konuyu zor öğrenmiş olmak, anlatırken karşımdakinin nerede ve neden takılacağını sezmemi kolaylaştırdı. Bu kitap, o yıllarda bocalarken eksikliğini derinden hissettiğim sağlam ve derli toplu zemini benden sonrakilere sunma çabasıdır.

Bu kitabı ben yazmadım; yapay zekâ araçlarından yararlanarak derledim ve düzenledim. Müfredat, öncelikler ve kapsam benim tecrübeme dayanıyor; yazım ve dizgi işlerinde de bu araçlardan yararlandım. Bu araçların ürettiği metnin belirlediğim kalite çizgisinin üstünde kalmasını denetledim ve gerekli düzeltmeleri yaptım.

Kitapla ilgili her türlü eleştiri ve geri bildiriminiz için bana \texttt{dseferoglu26@\allowbreak ku.edu.\allowbreak tr} adresinden ulaşabilirsiniz.

\clearpage